\documentclass[11pt,a4paper]{article}

% ============================================================
%  Informe: aprendizaje autosupervisado (SSL) en OASIS3
%  VICReg / BYOL / DINOv2 para clasificación AD/HC
% ============================================================

\usepackage[a4paper,margin=2.3cm]{geometry}
\usepackage[T1]{fontenc}
\usepackage{lmodern}
\usepackage[spanish]{babel}
\usepackage{microtype}
\usepackage{xcolor}
\usepackage{amsmath,amssymb}
\usepackage{booktabs}
\usepackage{array}
\usepackage{graphicx}
\usepackage{caption}
\usepackage{subcaption}
\usepackage{enumitem}
\usepackage{fancyhdr}
\usepackage{titlesec}
\usepackage{longtable}
\usepackage[most]{tcolorbox}
\usepackage{tikz}
\usetikzlibrary{arrows.meta,positioning,shapes.geometric,fit,backgrounds}

% ------------------------------------------------------------
% Paleta de colores
% ------------------------------------------------------------
\definecolor{primaryblue}{HTML}{1B4F72}
\definecolor{steelblue}{HTML}{2E86AB}
\definecolor{accentorange}{HTML}{E67E22}
\definecolor{hcgreen}{HTML}{1E8449}
\definecolor{adred}{HTML}{B03A2E}
\definecolor{softgray}{HTML}{85929E}
\definecolor{panelbg}{HTML}{EAF2F8}
\definecolor{warnbg}{HTML}{FDEBD0}
\definecolor{tldrbg}{HTML}{EBF5FB}

% ------------------------------------------------------------
% hyperref al final de la carga de paquetes
% ------------------------------------------------------------
\usepackage[colorlinks=true,linkcolor=primaryblue,citecolor=primaryblue,
            urlcolor=steelblue,linktoc=all]{hyperref}

% ------------------------------------------------------------
% Cajas de contenido
% ------------------------------------------------------------
\newtcolorbox{tldrbox}{colback=tldrbg,colframe=primaryblue,coltitle=white,
  colbacktitle=primaryblue,fonttitle=\bfseries,boxrule=0.9pt,arc=3pt,
  left=8pt,right=8pt,top=5pt,bottom=5pt,title=En una frase}
\newtcolorbox{ideabox}{colback=panelbg,colframe=steelblue,coltitle=white,
  colbacktitle=steelblue,fonttitle=\bfseries,boxrule=0.8pt,arc=2pt,
  left=8pt,right=8pt,top=5pt,bottom=5pt,title=Idea clave}
\newtcolorbox{warnbox}{colback=warnbg,colframe=accentorange,coltitle=white,
  colbacktitle=accentorange,fonttitle=\bfseries,boxrule=0.8pt,arc=2pt,
  left=8pt,right=8pt,top=5pt,bottom=5pt,title=Ojo con esto}

% ------------------------------------------------------------
% Encabezados y secciones
% ------------------------------------------------------------
\pagestyle{fancy}
\fancyhf{}
\renewcommand{\headrulewidth}{0.5pt}
\fancyhead[L]{\small\color{softgray}Informe SSL · OASIS3}
\fancyhead[R]{\small\color{softgray}VICReg / BYOL / DINOv2}
\fancyfoot[C]{\small\thepage}

\titleformat{\section}{\normalfont\Large\bfseries\color{primaryblue}}{\thesection}{0.6em}{}
\titleformat{\subsection}{\normalfont\large\bfseries\color{steelblue}}{\thesubsection}{0.6em}{}
\titleformat{\subsubsection}{\normalfont\bfseries\color{primaryblue!80!black}}{\thesubsubsection}{0.6em}{}

\graphicspath{{../campaign_full50/figures/}{../campaign_mri_fullview20/figures/}}

\setlength{\parskip}{4pt}
\setlist[itemize]{topsep=2pt,itemsep=2pt}
\setlist[enumerate]{topsep=2pt,itemsep=2pt}

% ------------------------------------------------------------
% Estilos TikZ comunes a todos los diagramas
% ------------------------------------------------------------
\tikzset{
  box/.style={rectangle, rounded corners=2pt, draw=primaryblue, thick,
    fill=primaryblue!6, text width=3.6cm, align=center, minimum height=1cm,
    inner sep=5pt, font=\small},
  wbox/.style={box, text width=7.2cm},
  splitbox/.style={rectangle, rounded corners=2pt, draw=hcgreen, thick,
    fill=hcgreen!8, text width=2.9cm, align=center, minimum height=1.1cm,
    inner sep=4pt, font=\scriptsize},
  viewbox/.style={rectangle, rounded corners=1pt, draw=softgray, thick,
    fill=white, text width=2.6cm, align=center, minimum height=0.9cm,
    inner sep=4pt, font=\scriptsize},
  onlinebox/.style={rectangle, rounded corners=2pt, draw=steelblue, thick,
    fill=steelblue!8, text width=3.4cm, align=center, minimum height=1cm,
    inner sep=4pt, font=\scriptsize},
  targetbox/.style={rectangle, rounded corners=2pt, draw=softgray, thick,
    dashed, fill=softgray!10, text width=3.4cm, align=center,
    minimum height=1cm, inner sep=4pt, font=\scriptsize},
  lossbox/.style={rectangle, rounded corners=2pt, draw=accentorange, thick,
    fill=accentorange!8, text width=8.4cm, align=left, inner sep=6pt,
    font=\scriptsize},
  notebox/.style={rectangle, rounded corners=2pt, draw=softgray, dashed,
    fill=softgray!6, text width=8.4cm, align=left, inner sep=5pt,
    font=\scriptsize\itshape},
  arr/.style={-{Latex[length=2.2mm]}, thick, primaryblue},
  arrg/.style={-{Latex[length=2mm]}, thin, softgray},
  emaarr/.style={-{Latex[length=2mm]}, thin, dashed, accentorange},
}

\begin{document}

% ============================================================
% PORTADA
% ============================================================
\begin{titlepage}
\centering
\vspace*{1.2cm}
{\Huge\bfseries\color{primaryblue} Aprender sin etiquetas\\[2pt] a partir de resonancias cerebrales}\\[10pt]
{\Large\color{steelblue} VICReg, BYOL y DINOv2 sobre OASIS3\\ para distinguir Alzheimer de controles sanos}\\[26pt]

\begin{tikzpicture}
\draw[primaryblue, thick] (0,0) -- (11,0);
\end{tikzpicture}

\vspace{18pt}
{\large \textit{Un recorrido visual por los datos, las tres arquitecturas\\
y lo que sus representaciones sí (y no) capturan sobre la enfermedad.}}

\vspace{40pt}
\begin{tcolorbox}[colback=panelbg,colframe=primaryblue,arc=4pt,width=0.82\textwidth,
  boxrule=1pt, left=14pt,right=14pt,top=10pt,bottom=10pt]
\begin{center}
\textbf{oasis3-ssl-workbench}\\[4pt]
\small 981 sujetos OASIS3 · 3 métodos de preentrenamiento autosupervisado\\
1 protocolo de evaluación común: un \emph{probe} lineal sobre encoder congelado
\end{center}
\end{tcolorbox}

\vfill
{\small Informe generado a partir de la campaña \texttt{campaign\_full50} (981 sujetos,
hasta 50 épocas) y de la sesión de ajuste del 20 de septiembre de 2026.}\\[4pt]
{\small 20 de septiembre de 2026}
\end{titlepage}

\tableofcontents
\newpage

% ============================================================
\section{Por qué este informe}
% ============================================================

Este documento explica, de la forma más visual posible, tres cosas:

\begin{enumerate}
\item cómo un conjunto de resonancias magnéticas cerebrales (OASIS3) se convierte
  en los tensores que consume una red neuronal;
\item qué le pasa exactamente a \emph{una} resonancia de \emph{un} paciente al
  entrar a cada uno de los tres modelos autosupervisados del proyecto —VICReg,
  BYOL y DINOv2—; y
\item qué tan bien, al final, esas representaciones aprendidas \emph{sin
  etiquetas} sirven para distinguir personas con Alzheimer (AD) de controles
  sanos (HC) cuando se les pega encima un clasificador simple.
\end{enumerate}

\begin{tldrbox}
Entrenamos tres redes a "mirar" cortes axiales de resonancias cerebrales sin
decirles nunca si la persona tiene Alzheimer. Después las congelamos y les
pedimos, con un clasificador lineal muy simple, que separen AD de HC. VICReg
es el único que mejora de forma consistente sobre un encoder sin entrenar;
BYOL y DINOv2 quedan dentro del margen de ruido en la campaña oficial de 50
épocas, aunque una sesión posterior con otros backbones movió bastante esas
cifras (Sección~\ref{sec:extension}).
\end{tldrbox}

Todo el código, configuraciones y figuras de este informe viven en el
repositorio \texttt{oasis3-ssl-workbench}; al final (Sección~\ref{sec:repro})
se listan los comandos exactos para reproducir cada número.

% ============================================================
\section{El dataset: de un sujeto clínico a un tensor 2D}
\label{sec:data}
% ============================================================

OASIS3 aporta volúmenes de resonancia magnética estructural (T1) registrados
a un espacio estándar (MNI), junto con metadatos clínicos por sesión: sujeto,
edad, y grupo diagnóstico (entre otros, \texttt{AD}, \texttt{HC}, y
diagnósticos intermedios que aquí tratamos como ``otro''). Un mismo sujeto
puede tener \emph{varias} sesiones (visitas) a lo largo de los años.

\subsection{El recorrido completo, de arriba a abajo}

\begin{figure}[htbp]
\centering
\begin{tikzpicture}[node distance=6mm]

\node[wbox] (csv) at (0,15) {\textbf{CSV clínico OASIS3}\\981 sujetos · columna \texttt{group} (AD / HC / otro)};
\node[wbox] (lmdb) at (8.6,15) {\textbf{LMDB de volúmenes MNI}\\intensidad \texttt{int16}, ya registrados};

\node[wbox, fill=hcgreen!8, draw=hcgreen] (split) at (4.3,12.6) {\textbf{Split fijo por sujeto}\\un sujeto vive en un único conjunto: nunca hay fuga entre train/val/test};

\node[splitbox] (train) at (0.3,10.1) {\textbf{Train}\\685 sujetos\\1\,226 sesiones disponibles};
\node[splitbox] (val)   at (4.3,10.1) {\textbf{Val}\\146 sujetos\\1 sesión fija c/u};
\node[splitbox] (test)  at (8.3,10.1) {\textbf{Test}\\150 sujetos\\1 sesión fija c/u};

\node[wbox, text width=9.5cm] (slice) at (4.3,6.7) {\textbf{Elegir sesión y corte axial}\\[2pt]train: sesión y fracción aleatorias en $[0.30,0.70]$ del eje $z$ (cambian cada época)\\[2pt]val/test: sesión y fracción fijas ($0.50$, el corte central)};

\node[wbox, text width=9.5cm] (norm) at (4.3,3.9) {\textbf{Normalización por percentiles}\\$\rightarrow$ tensor $1\times H\times W$ en $[0,1]$};

\node[onlinebox, text width=3.6cm] (augtr) at (1.2,1.2) {\texttt{SSLViewTransform}\\crop aleatorio + flips + rotación + blur + ruido\\$\Rightarrow$ \textbf{2 vistas} independientes (solo \emph{train})};
\node[targetbox, text width=3.6cm] (augev) at (7.4,1.2) {\texttt{EvalTransform}\\solo \emph{resize} determinista\\$\Rightarrow$ 1 vista reproducible (val/test/\emph{probe})};

\node[box, fill=accentorange!10, draw=accentorange, text width=7.6cm] (enc) at (4.3,-1.6)
  {Entra al \textbf{encoder} de VICReg, BYOL o DINOv2\\(ver Sección~\ref{sec:models})};

\draw[arr] (csv) -- (split);
\draw[arr] (lmdb) -- (split);
\draw[arr] (split) -- (train);
\draw[arr] (split) -- (val);
\draw[arr] (split) -- (test);
\draw[arr] (train) -- (slice);
\draw[arr] (val) -- (slice);
\draw[arr] (test) -- (slice);
\draw[arr] (slice) -- (norm);
\draw[arr] (norm) -- (augtr);
\draw[arr] (norm) -- (augev);
\draw[arr] (augtr) -- (enc);
\draw[arr] (augev) -- (enc);
\end{tikzpicture}
\caption{De 981 sujetos clínicos a un tensor 2D listo para el encoder. El
split por sujeto se decide una sola vez y es el mismo para las tres
arquitecturas.}
\label{fig:pipeline}
\end{figure}

\begin{ideabox}
El \emph{split} nunca mezcla sesiones de un mismo sujeto entre conjuntos. Si
no fuera así, el modelo podría ``reconocer'' el cerebro de una persona que ya
vio en \emph{train} al evaluarlo en \emph{test}, e inflar artificialmente el
desempeño. El código (\texttt{select\_rows} en
\texttt{oasis3\_lmdb\_datamodule.py}) verifica activamente que los conjuntos
de sujetos no se solapen y lanza un error si detecta fuga.
\end{ideabox}

\subsection{Una sesión, un corte, dos miradas}

Cada volumen 3D se recorta a un único corte axial 2D (plano perpendicular al
eje cráneo-pies). Durante el preentrenamiento SSL, \emph{train} muestrea en
cada época una sesión al azar del sujeto y una fracción aleatoria del eje $z$
dentro de $[0.30, 0.70]$ (evita cortes demasiado periféricos); \emph{val} y
\emph{test} siempre usan la misma sesión ``primaria'' y el corte central
($0.50$), para que la evaluación sea reproducible.

A partir de ese corte, \texttt{SSLViewTransform} genera \textbf{dos vistas
aumentadas independientes} de la misma imagen: recorte aleatorio y reescalado
(\emph{scale} $\in[0.3,1.0]$), flips horizontal/vertical, rotación
($\pm15^\circ$), \emph{blur} gaussiano ocasional y ruido gaussiano de baja
intensidad. Ninguna de las tres arquitecturas ve nunca la etiqueta AD/HC en
este paso: la única señal es ``estas dos imágenes vienen del mismo cerebro''.

\begin{figure}[htbp]
\centering
\includegraphics[width=0.92\linewidth]{augmentation_examples.png}
\caption{Dos políticas de aumentación aplicadas al mismo corte axial de
entrenamiento, generando dos vistas por política (ilustración tomada de la
campaña \texttt{mri\_fullview20}, que comparte el mismo código de
aumentación; la campaña de 50 épocas usa una configuración de recorte más
agresiva, \texttt{scale}$\in[0.3,1.0]$, ver Tabla~\ref{tab:overview}).}
\label{fig:augmentations}
\end{figure}

DINOv2 añade una variante propia (\texttt{DINOv2ViewTransform}): además de las
dos vistas ``globales'' de $120\times120$, genera varios recortes
\textbf{locales} más pequeños ($60\times60$, tomados de una región mucho más
acotada del cerebro, \emph{scale}$\in[0.1,0.4]$). La idea, típica de la
familia DINO, es forzar a que una vista pequeña y local sea consistente con la
representación global del mismo cerebro.

\subsection{Cuántos sujetos hay realmente en cada paso}

No todos los sujetos con imagen tienen una etiqueta AD/HC inequívoca (algunos
tienen diagnósticos intermedios, o etiquetas que cambian entre sesiones). El
preentrenamiento SSL usa \textbf{todos} los sujetos del split, con o sin
etiqueta; el \emph{probe} lineal posterior (Sección~\ref{sec:probe}) sólo
puede usar los que sí tienen una etiqueta AD/HC clara.

\begin{table}[htbp]
\centering
\small
\begin{tabular}{lccc}
\toprule
 & \textbf{Train} & \textbf{Val} & \textbf{Test} \\
\midrule
Sujetos totales en el split (SSL) & 685 & 146 & 150 \\
Sesiones disponibles para SSL & 1\,226 & 146 & 150 \\
Sujetos AD/HC inequívocos (\emph{probe}) & 615 & 131 & 134 \\
\quad de los cuales AD & --- & --- & 31 \\
\quad de los cuales HC & --- & --- & 103 \\
\bottomrule
\end{tabular}
\caption{Sujetos disponibles para preentrenamiento SSL frente a los
realmente usados por el \emph{probe} AD/HC, campaña de 50 épocas.}
\label{tab:subjects}
\end{table}

% ============================================================
\section{Tres formas de aprender sin etiquetas}
\label{sec:models}
% ============================================================

Las tres arquitecturas comparten la misma idea de fondo —tomar dos vistas
aumentadas de la misma imagen y enseñarle a la red que ``vienen de lo
mismo''— pero difieren en \emph{cómo} evitan la trampa obvia: si sólo
premiamos que las dos vistas se parezcan, la red puede colapsar a una
constante y ``ganar'' trivialmente. Cada método tiene su propio antídoto
contra ese colapso.

\subsection{El ingrediente común: el encoder}

En los tres casos, el \emph{encoder} (o \emph{backbone}) es la pieza que
después se congela y se reutiliza para clasificar AD/HC —es lo único que nos
importa conservar del preentrenamiento SSL. Recibe una imagen de un canal
($1\times H\times W$, un corte axial) y entrega un vector de tamaño fijo.

\begin{table}[htbp]
\centering
\small
\begin{tabular}{lccc}
\toprule
 & \textbf{VICReg} & \textbf{BYOL} & \textbf{DINOv2} \\
\midrule
Backbone & ConvNeXt-tiny & ConvNeXt-tiny & ViT-Small (\emph{patch} 60) \\
Preentrenado en ImageNet & no & no & no \\
Dimensión de salida & 768 & 768 & 384 \\
Tamaño de imagen & $96\times96$ & $96\times96$ & $120\times120$ \\
Vistas por paso & 2 & 2 & 2 globales + 2 locales \\
\bottomrule
\end{tabular}
\caption{Backbones usados en la campaña oficial de 50 épocas
(\texttt{campaign\_full50}). La sesión exploratoria del 20-09 cambió algunos
de estos backbones; ver Sección~\ref{sec:extension}.}
\label{tab:backbones}
\end{table}

\subsection{VICReg: parecerse sin colapsar, a punta de estadística}

\textbf{Idea:} una sola red procesa las dos vistas (no hay red ``espejo'' ni
promedio móvil). Para evitar el colapso, no basta con que las dos
proyecciones se parezcan (\emph{invarianza}): además se exige, \emph{sobre el
lote}, que cada dimensión del embedding tenga varianza mínima
(\emph{varianza}) y que las dimensiones no queden correlacionadas entre sí
(\emph{covarianza}). Son tres términos que tironean en direcciones
distintas: parecerse, pero sin aplanarse, y sin repetir la misma información
en varias dimensiones.

\begin{figure}[htbp]
\centering
\begin{tikzpicture}
\node[viewbox] (v1) at (1.2,10.4) {Vista 1};
\node[viewbox] (v2) at (9.2,10.4) {Vista 2};

\node[box] (e1) at (1.2,8.6) {Encoder $f_\theta$\\ConvNeXt-tiny\\$h\in\mathbb{R}^{768}$};
\node[box] (e2) at (9.2,8.6) {Encoder $f_\theta$\\ConvNeXt-tiny\\$h\in\mathbb{R}^{768}$};
\draw[dashed, softgray] (e1) -- (e2) node[midway, above, font=\tiny, text=softgray] {mismos pesos $\theta$ (no hay red espejo)};

\node[box] (p1) at (1.2,6.6) {Proyector $g_\theta$\\MLP 3 capas, BN+ReLU\\$768\to2048\to2048$};
\node[box] (p2) at (9.2,6.6) {Proyector $g_\theta$\\MLP 3 capas, BN+ReLU\\$768\to2048\to2048$};
\draw[dashed, softgray] (p1) -- (p2);

\node[font=\small] (z1) at (1.2,5.2) {$z_1\in\mathbb{R}^{2048}$};
\node[font=\small] (z2) at (9.2,5.2) {$z_2\in\mathbb{R}^{2048}$};

\node[lossbox] (loss) at (5.2,2.9) {
  \textbf{Pérdida VICReg} $= 25\cdot\underbrace{\mathrm{MSE}(z_1,z_2)}_{\text{invarianza}}
  \;+\; 25\cdot\underbrace{\textstyle\sum_j\max(0,\,1-\sigma_j)}_{\text{varianza, por dimensión}}
  \;+\; 1\cdot\underbrace{\textstyle\sum_{i\ne j} C_{ij}^2}_{\text{covarianza, fuera de la diagonal}}$
};

\node[notebox] (note) at (5.2,0.6) {Sin red EMA, sin pares negativos. Opcionalmente existe un banco (\emph{queue}) de proyecciones pasadas para estabilizar varianza/covarianza en lotes chicos --- desactivado (\texttt{queue\_size=0}) en estas corridas.};

\draw[arr] (v1) -- (e1);
\draw[arr] (v2) -- (e2);
\draw[arr] (e1) -- (p1);
\draw[arr] (e2) -- (p2);
\draw[arr] (p1) -- (z1);
\draw[arr] (p2) -- (z2);
\draw[arr] (z1) -- (loss.north west);
\draw[arr] (z2) -- (loss.north east);
\end{tikzpicture}
\caption{Un paciente atravesando VICReg: la \emph{misma} red ve las dos
vistas; el colapso se evita con los términos de varianza y covarianza, no
con una arquitectura asimétrica.}
\label{fig:vicreg}
\end{figure}

\begin{ideabox}
En VICReg no hay ``profesor'' ni ``alumno'': una única red procesa ambas
vistas en el mismo paso hacia adelante. Toda la lucha contra el colapso pasa
por la forma de la pérdida, no por la arquitectura.
\end{ideabox}

\subsection{BYOL: perseguir un reflejo que se mueve despacio}

\textbf{Idea:} en vez de una pérdida estadística, BYOL usa dos redes: una
\emph{online} (la que aprende por gradiente) y una \emph{target} (una copia
cuyos pesos son un promedio móvil exponencial ---EMA--- de la \texttt{online},
y que \textbf{nunca} recibe gradientes). La red online, además, tiene un
\emph{predictor} extra que intenta anticipar la proyección que la red target
produciría para la \emph{otra} vista. Como el objetivo (target) se mueve
mucho más despacio que el alumno, y nunca es exactamente igual, colapsar a
una constante deja de ser una solución trivial: el predictor tendría que
``adivinar'' un blanco que además se sigue moviendo.

\begin{figure}[htbp]
\centering
\begin{tikzpicture}
\node[viewbox] (v1) at (1.2,11.6) {Vista 1};
\node[viewbox] (v2) at (9.2,11.6) {Vista 2};

\node[onlinebox] (eo1) at (1.2,9.8) {Encoder \emph{online} $f_\theta$\\ConvNeXt-tiny};
\node[onlinebox] (eo2) at (9.2,9.8) {Encoder \emph{online} $f_\theta$\\ConvNeXt-tiny};
\draw[dashed, softgray] (eo1) -- (eo2) node[midway, above, font=\tiny, text=softgray] {mismos pesos $\theta$};

\node[onlinebox] (po1) at (1.2,8.1) {Proyector $g_\theta$\\$768\to512\to256$};
\node[onlinebox] (po2) at (9.2,8.1) {Proyector $g_\theta$\\$768\to512\to256$};

\node[onlinebox] (pr1) at (1.2,6.4) {Predictor $q_\theta$\\256 (2 capas)\\$\Rightarrow p_1$};
\node[onlinebox] (pr2) at (9.2,6.4) {Predictor $q_\theta$\\256 (2 capas)\\$\Rightarrow p_2$};

\node[targetbox] (et1) at (1.2,4.0) {Encoder \emph{target} $f_\xi$ (EMA)\\\texttt{sin gradiente}};
\node[targetbox] (et2) at (9.2,4.0) {Encoder \emph{target} $f_\xi$ (EMA)\\\texttt{sin gradiente}};

\node[targetbox] (pt1) at (1.2,2.3) {Proyector $g_\xi$ (EMA)\\$\Rightarrow t_1$};
\node[targetbox] (pt2) at (9.2,2.3) {Proyector $g_\xi$ (EMA)\\$\Rightarrow t_2$};

\node[lossbox, text width=9.6cm] (loss) at (5.2,0.2)
  {\textbf{Pérdida BYOL (simétrica)} $=\tfrac12\big[(2-2\cos(p_1,t_2)) + (2-2\cos(p_2,t_1))\big]$
  --- cada predicción se compara contra el target de la \emph{otra} vista; $t_1,t_2$ nunca reciben gradiente.};

\draw[arr] (v1) -- (eo1); \draw[arr] (v2) -- (eo2);
\draw[arr] (eo1) -- (po1); \draw[arr] (eo2) -- (po2);
\draw[arr] (po1) -- (pr1); \draw[arr] (po2) -- (pr2);
\draw[arrg] (v1) to[bend right=15] (et1); \draw[arrg] (v2) to[bend left=15] (et2);
\draw[arrg] (et1) -- (pt1); \draw[arrg] (et2) -- (pt2);
\draw[emaarr] (eo1.south west) to[bend right=35] node[midway,left,font=\tiny,text=accentorange]{EMA} (et1.north west);
\draw[emaarr] (eo2.south east) to[bend left=35] node[midway,right,font=\tiny,text=accentorange]{EMA} (et2.north east);
\draw[arr] (pr1) -- (loss.north west);
\draw[arr] (pr2) -- (loss.north east);
\draw[arrg] (pt1) -- (loss.west);
\draw[arrg] (pt2) -- (loss.east);
\end{tikzpicture}
\caption{Un paciente atravesando BYOL: la rama azul (\emph{online}) aprende
por gradiente; la rama gris punteada (\emph{target}) es una copia EMA que
sólo se actualiza por interpolación, nunca por \emph{backprop}. La
actualización EMA usa un coeficiente de momento que sube con \emph{coseno}
desde $0.996$ hacia $1$ a lo largo del entrenamiento.}
\label{fig:byol}
\end{figure}

\begin{warnbox}
No confundir el ``predictor'' de BYOL con un clasificador de AD/HC: sólo
predice la proyección de la \emph{otra} vista aumentada de la misma imagen.
La etiqueta clínica sigue sin aparecer en ningún punto de este diagrama.
\end{warnbox}

\subsection{DINOv2: un profesor sereno y un alumno que trabaja a ciegas}

\textbf{Idea:} de nuevo hay una red \emph{estudiante} (con gradiente) y una
\emph{profesora} EMA (sin gradiente), como en BYOL, pero aquí ambas son
\emph{Vision Transformers} y compiten sobre \emph{tokens}, no sobre un solo
vector. El estudiante ve las vistas globales con una fracción de sus parches
\textbf{enmascarados} (ocultos) y también recortes locales mucho más
pequeños; la profesora siempre ve las vistas globales completas, sin máscara.
La pérdida combina tres términos:

\begin{itemize}
\item \textbf{Distilación CLS cruzada} (``dino''): el token de clase (CLS)
  de cada vista del estudiante debe imitar la distribución de probabilidad
  que la profesora produce para una vista global \emph{distinta}.
\item \textbf{iBOT}: el estudiante también debe adivinar, a partir del
  contexto, qué distribución de parche predijo la profesora exactamente en
  los parches que el estudiante tuvo \emph{ocultos}.
\item \textbf{KoLeo}: un término de dispersión que empuja a los embeddings
  del lote a no amontonarse todos cerca unos de otros.
\end{itemize}

La profesora no se vuelve un blanco trivial (todo unos, por ejemplo) porque
sus logits se \emph{centran} (se les resta un promedio móvil de activaciones
pasadas) y se afinan con una temperatura baja antes del \emph{softmax}; esa
temperatura, además, sube gradualmente ($0.01\to0.02$ en las primeras 15
épocas) para no sobre-confiar en un profesor que todavía no sabe nada.

\begin{figure}[htbp]
\centering
\begin{tikzpicture}
\node[viewbox] (g1) at (1.4,11.2) {Recorte global 1\\$120\times120$};
\node[viewbox] (g2) at (5.4,11.2) {Recorte global 2\\$120\times120$};
\node[viewbox] (loc) at (9.4,11.2) {$n$ recortes locales\\$60\times60$ (p.ej.\ 2)};

\node[onlinebox, text width=9.4cm] (stud) at (5.4,9.2)
 {\textbf{Estudiante} --- ViT-Small, \emph{patch} 60, 384-d\\
 globales: 30\,\% de los parches \textbf{enmascarados}\ · locales: sin máscara, solo CLS};

\node[onlinebox] (clss) at (2.6,7.1) {Proyector CLS $g_\theta$\\$384\to512\to1024$};
\node[onlinebox] (patchs) at (8.2,7.1) {Cabezal de parches\\(solo parches enmascarados)};

\node[targetbox, text width=9.4cm] (teach) at (5.4,5.0)
 {\textbf{Profesora (EMA)} --- ViT-Small, \emph{patch} 60\\
 ve los 2 recortes globales completos, \texttt{sin máscara}, \texttt{sin gradiente}};

\node[targetbox] (clst) at (2.6,3.0) {Proyector CLS $g_\xi$};
\node[targetbox] (patcht) at (8.2,3.0) {Cabezal de parches $\xi$};

\node[targetbox, text width=8cm] (center) at (5.4,1.1)
 {Centrado (EMA de medias) + \emph{softmax} con $\tau_{\text{prof}}$ baja ($0.01\to0.02$, \emph{warm-up} 15 épocas)};

\node[lossbox, text width=10cm] (loss) at (5.4,-1.3)
 {\textbf{Pérdida DINOv2} $=$ distilación CLS cruzada (todas las vistas del
 estudiante contra las 2 vistas globales de la profesora, excluyendo la
 propia) $+\ 1.0\cdot$iBOT (parches enmascarados) $+\ 0.1\cdot$KoLeo
 (dispersión)};

\draw[arr] (g1) -- (stud); \draw[arr] (g2) -- (stud); \draw[arr] (loc) -- (stud);
\draw[arr] (stud) -- (clss); \draw[arr] (stud) -- (patchs);
\draw[arrg] (g1) to[bend right=25] (teach); \draw[arrg] (g2) to[bend left=25] (teach);
\draw[arrg] (teach) -- (clst); \draw[arrg] (teach) -- (patcht);
\draw[arrg] (clst) -- (center); \draw[arrg] (patcht) -- (center);
\draw[emaarr] (stud.south) to[bend right=20] node[midway,left,font=\tiny,text=accentorange]{EMA} (teach.north);
\draw[arr] (clss) -- (loss.north west);
\draw[arr] (patchs) -- (loss.north east);
\draw[arrg] (center) -- (loss.north);
\end{tikzpicture}
\caption{Un paciente atravesando DINOv2: el estudiante trabaja con
información incompleta (parches ocultos, recortes locales pequeños) y debe
igualar, de todos modos, lo que la profesora --- que sí ve todo, pero nunca
aprende por gradiente --- observó en las vistas completas.}
\label{fig:dinov2}
\end{figure}

\subsection{Las tres arquitecturas, una tabla}

\begin{table}[htbp]
\centering
\small
\begin{tabular}{p{2.1cm}p{3.4cm}p{3.4cm}p{3.6cm}}
\toprule
 & \textbf{VICReg} & \textbf{BYOL} & \textbf{DINOv2} \\
\midrule
Nº de redes & 1 & 2 (\emph{online} + EMA) & 2 (estudiante + EMA) \\
Anti-colapso & varianza + covarianza (estadística) & red target EMA + predictor & \emph{masking} + centrado + KoLeo \\
Usa negativos & no & no & no \\
Vistas & 2 globales & 2 globales & 2 globales + locales \\
Pesos SSL & $\lambda_{inv}{=}25$, $\lambda_{var}{=}25$, $\lambda_{cov}{=}1$ & EMA base $0.996\to1$ (coseno) & EMA base $0.996\to1$, iBOT $=1.0$, KoLeo $=0.1$ \\
Riesgo típico de colapso & varianza $\to0$ en todas las dimensiones & representación constante si el predictor ``hace trampa'' & profesora sobre-confiada desde el inicio \\
\bottomrule
\end{tabular}
\caption{Comparación estructural de los tres métodos, campaña oficial de 50 épocas.}
\label{tab:overview}
\end{table}

% ============================================================
\section{Cómo se evalúa: el \emph{probe} lineal sobre encoder congelado}
\label{sec:probe}
% ============================================================

Preentrenar sin etiquetas es solo la mitad de la historia: para saber si
sirvió, se \textbf{congela} el encoder (sin más ajuste de sus pesos) y se
entrena únicamente una regresión logística encima. Este protocolo --- el
\emph{frozen linear probe} --- es el estándar en la literatura SSL
precisamente porque aísla qué tan \emph{linealmente separable} quedó la
información clínicamente relevante en el espacio de embeddings, sin dejar
que un clasificador no lineal potente disimule una mala representación.

\begin{figure}[htbp]
\centering
\begin{tikzpicture}
\node[onlinebox] (enc) at (2.6,10.4) {Encoder congelado\\\textbf{pesos SSL entrenados}};
\node[targetbox] (encr) at (9.4,10.4) {Encoder congelado\\\textbf{pesos aleatorios} (control, misma arquitectura sin entrenar)};

\node[box, text width=3.6cm] (emb) at (2.6,8.5) {Embedding por sesión evaluada\\(1 vista determinista)};
\node[box, text width=3.6cm, draw=softgray] (embr) at (9.4,8.5) {Embedding por sesión evaluada\\(1 vista determinista)};

\node[box, text width=3.6cm] (avg) at (2.6,6.6) {Promedio por sujeto\\(si tiene varias sesiones evaluadas)};
\node[box, text width=3.6cm, draw=softgray] (avgr) at (9.4,6.6) {Promedio por sujeto};

\node[box, text width=3.6cm] (fit) at (2.6,4.7) {Ajustar en \emph{train}\\615 sujetos AD/HC\\\texttt{StandardScaler+LogReg}};
\node[box, text width=3.6cm, draw=softgray] (fitr) at (9.4,4.7) {Ajustar en \emph{train}\\615 sujetos AD/HC};

\node[box, text width=3.6cm] (sel) at (2.6,2.8) {Elegir $C\in\{0.01,0.1,1,10\}$\\por AUROC en \emph{val}\\131 sujetos};
\node[box, text width=3.6cm, draw=softgray] (selr) at (9.4,2.8) {Elegir $C$\\por AUROC en \emph{val}};

\node[box, fill=accentorange!10, draw=accentorange, text width=3.6cm] (test) at (2.6,0.9) {AUROC en \emph{test}\\134 sujetos\\(31 AD / 103 HC)};
\node[box, fill=accentorange!10, draw=accentorange, text width=3.6cm] (testr) at (9.4,0.9) {AUROC en \emph{test}\\misma cohorte};

\node[lossbox, text width=9.6cm] (comp) at (6,-1.3)
 {\textbf{Bootstrap pareado por sujeto} (2\,000 remuestreos)\\
 $\Delta\text{AUROC} = \text{AUROC}_{\text{SSL}} - \text{AUROC}_{\text{aleatorio}}$ \ $\Rightarrow$ IC 95\,\%};

\draw[arr] (enc) -- (emb); \draw[arrg] (encr) -- (embr);
\draw[arr] (emb) -- (avg); \draw[arrg] (embr) -- (avgr);
\draw[arr] (avg) -- (fit); \draw[arrg] (avgr) -- (fitr);
\draw[arr] (fit) -- (sel); \draw[arrg] (fitr) -- (selr);
\draw[arr] (sel) -- (test); \draw[arrg] (selr) -- (testr);
\draw[arr] (test) -- (comp.north west);
\draw[arr] (testr) -- (comp.north east);
\end{tikzpicture}
\caption{Protocolo del \emph{probe} lineal: cada encoder (SSL vs.\ control
aleatorio) recorre exactamente el mismo camino, con el mismo split y la
misma cohorte de sujetos, para que la única diferencia posible sea la
calidad de la representación.}
\label{fig:probe}
\end{figure}

\begin{ideabox}
El control ``aleatorio'' no es un truco menor: un encoder \emph{sin
entrenar} ya tiene, gracias a su arquitectura, cierta capacidad de separar
clases cuando se le ajusta encima un clasificador. Comparar contra ese
control --- y no contra el azar puro (AUROC $=0.5$) --- es lo que permite
afirmar si el preentrenamiento SSL agregó algo por encima de lo que la
arquitectura ya regalaba gratis.
\end{ideabox}

Además del AUROC, el script calcula el \textbf{rango efectivo espectral} de
los embeddings de \emph{train} (cuántas direcciones del espacio realmente se
usan, en un sentido continuo) y la \textbf{desviación estándar media por
dimensión}: dos formas de detectar colapso parcial que un AUROC por sí solo
no muestra.

% ============================================================
\section{Resultados: la campaña oficial de 50 épocas}
\label{sec:results}
% ============================================================

Protocolo: 981 sujetos OASIS3, split fijo y disjunto por sujeto
(685/146/150), configuraciones nativas de cada método (sin ajustar
hiperparámetros mirando las etiquetas), hasta 50 épocas con \emph{early
stopping} basado únicamente en \texttt{val/loss} SSL (paciencia 10, mejora
mínima $0.001$) para los modelos que lo activaron.

\subsection{Curvas de pérdida SSL}

\begin{figure}[htbp]
\centering
\includegraphics[width=0.72\linewidth]{ssl_learning_curves.png}
\caption{Pérdida de entrenamiento y validación por método. Las escalas
\textbf{no} son comparables entre métodos (cada pérdida se define distinto);
lo relevante es la forma de cada curva. VICReg y BYOL activaron
\emph{early stopping} en la época 36 (mejor validación en la época 26);
DINOv2 completó las 50 épocas sin activarlo.}
\label{fig:curves}
\end{figure}

\begin{warnbox}
La curva de validación de BYOL es muy ruidosa (sube y baja bruscamente entre
épocas) y la de VICReg apenas se mueve tras las primeras iteraciones: ambas
son señales de que, aunque la pérdida SSL \emph{funciona} como objetivo de
optimización, no hay garantía de que decrezca monótonamente ni de que un
menor valor final implique una mejor representación para AD/HC.
\end{warnbox}

\subsection{¿El encoder entrenado ve más que uno al azar?}

\begin{table}[htbp]
\centering
\small
\begin{tabular}{lccccc}
\toprule
\textbf{Método} & \textbf{Épocas} & \textbf{AUROC SSL} & \textbf{AUROC aleatorio} & \boldmath$\Delta$\textbf{AUROC [IC 95\,\%]} \\
\midrule
DINOv2 & 50 (última época) & 0.717 & 0.720 & $-0.003$ [$-0.065$, $+0.056$] \\
VICReg & 36 (mejor val.\ SSL) & \textbf{0.776} & 0.748 & $\mathbf{+0.028}$ [$-0.029$, $+0.089$] \\
BYOL & 36 (mejor val.\ SSL) & 0.736 & 0.748 & $-0.012$ [$-0.068$, $+0.042$] \\
\bottomrule
\end{tabular}
\caption{AUROC del \emph{probe} AD/HC por sujeto en \emph{test} (134
sujetos: 31 AD, 103 HC). Intervalos de confianza por \emph{bootstrap}
pareado por sujeto, condicionados a un único split y semilla.}
\label{tab:results}
\end{table}

\begin{figure}[htbp]
\centering
\includegraphics[width=0.9\linewidth]{probe_vs_baseline_table.png}
\caption{La misma tabla, en formato visual.}
\end{figure}

\begin{figure}[htbp]
\centering
\begin{subfigure}[t]{0.66\linewidth}
\includegraphics[width=\linewidth]{ad_hc_subject_auc.png}
\caption{AUROC absoluto: encoder SSL vs.\ control aleatorio, por método.}
\end{subfigure}
\hfill
\begin{subfigure}[t]{0.30\linewidth}
\includegraphics[width=\linewidth]{ad_hc_delta_auc.png}
\caption{$\Delta$AUROC (SSL menos aleatorio) con IC 95\,\%. Solo VICReg
queda del lado positivo, y con IC amplio.}
\end{subfigure}
\caption{Ninguno de los tres intervalos de confianza excluye claramente el
cero: con 134 sujetos de test, esta campaña no alcanza para declarar una
mejora \emph{estadísticamente} contundente, aunque VICReg es el único que
apunta consistentemente en la dirección correcta.}
\label{fig:auc}
\end{figure}

\subsection{¿Calibración, o solo ranking?}

\begin{figure}[htbp]
\centering
\includegraphics[width=0.7\linewidth]{ad_hc_probe_log_loss.png}
\caption{\emph{Log-loss} en test. El \emph{prior} (predecir siempre la
prevalencia de AD del set de \emph{train}, sin mirar la imagen) logra el
menor \emph{log-loss} de los cuatro; los tres \emph{probes} usan
ponderación de clases y eligen $C$ por AUROC, no por calibración, así que
un buen ranking (AUROC) no implica probabilidades bien calibradas.}
\label{fig:logloss}
\end{figure}

\subsection{Dispersión de los embeddings: ¿colapso parcial?}

\begin{table}[htbp]
\centering
\small
\begin{tabular}{l*{4}{>{\centering\arraybackslash}p{2.35cm}}}
\toprule
\textbf{Método} & \textbf{Rango efectivo SSL} & \textbf{Rango efectivo aleatorio} & \textbf{Desv.\ media SSL} & \textbf{Desv.\ media aleatoria} \\
\midrule
DINOv2 & 2.0 & 7.8 & 0.0571 & 0.0237 \\
VICReg & 47.2 & 105.8 & 0.0658 & 0.0241 \\
BYOL & 2.5 & 105.8 & 0.0532 & 0.0241 \\
\bottomrule
\end{tabular}
\caption{Rango efectivo espectral de los embeddings de \emph{train}
(cuántas direcciones realmente se usan) y desviación estándar media por
dimensión.}
\label{tab:rank}
\end{table}

\begin{warnbox}
BYOL y DINOv2 concentran casi toda su varianza en un puñado de direcciones
(rango efectivo $\approx2$--$2.5$, frente a $\approx8$--$106$ del control
aleatorio): sus representaciones sí sufren \textbf{colapso parcial}, y no
superan al control en AUROC. VICReg conserva muchas más direcciones activas
(47.2), y es el único que mejora sobre su control --- aunque el intervalo de
confianza de esa mejora todavía incluye el cero. Un rango efectivo bajo no es
un defecto teórico abstracto: es la explicación mecánica más probable de por
qué BYOL y DINOv2 no logran superar a un encoder sin entrenar en esta
campaña.
\end{warnbox}

\subsection{Síntesis de la campaña oficial}

\begin{itemize}
\item \textbf{VICReg}: mejora modesta pero consistente (+0.028 AUROC) y es
  el único método que conserva una representación con buena dispersión
  (rango efectivo 47.2). Evidencia estadística todavía débil.
\item \textbf{BYOL}: sin mejora sobre el control aleatorio en esta
  configuración (ConvNeXt-tiny sin preentrenar); su representación muestra
  colapso parcial marcado.
\item \textbf{DINOv2}: esencialmente empatado con el azar; su
  representación es la más colapsada de las tres (rango efectivo 2.0).
\end{itemize}

% ============================================================
\section{Extensión exploratoria: sesión del 20 de septiembre de 2026}
\label{sec:extension}
% ============================================================

Después de la campaña oficial, se probó si cambiar el \emph{backbone} de
VICReg y BYOL --- usando pesos preentrenados en ImageNet en vez de
inicialización aleatoria --- movía la aguja para OASIS3. Esta sección resume
esos hallazgos con las salvedades necesarias: \textbf{no están todavía
persistidos en el repositorio} (viven en \texttt{logs/} y en un directorio
temporal del sistema) y arrastran un problema metodológico identificado
durante la propia sesión.

\subsection{Qué cambió respecto a la campaña oficial}

\begin{itemize}
\item \texttt{MLPProjector} ganó un parámetro \texttt{dropout} (por defecto
  $0$, no afecta a BYOL ni DINOv2, que no lo usan).
\item VICReg pasó a usar ConvNeXt-tiny \textbf{preentrenado en ImageNet}
  (\texttt{model.backbone.\allowbreak pretrained=true}) y, tras detectar
  sobreajuste, un \texttt{dropout=0.1} en el proyector.
\item BYOL pasó de ConvNeXt-tiny (config oficial,
  \texttt{byol\_convnext.yaml}) a \textbf{ResNet-18 preentrenado en
  ImageNet} (\texttt{byol\_resnet18.yaml}), replicando la configuración de
  un informe previo del curso EL7006.
\item DINOv2 mantuvo la misma arquitectura (ViT-Small, \emph{patch} 60) y se
  reentrenó como control de estabilidad entre sesiones.
\end{itemize}

\subsection{Resultados preliminares}

\begin{table}[htbp]
\centering
\small
\begin{tabular}{p{2.9cm}p{3.0cm}cp{4.6cm}}
\toprule
\textbf{Modelo} & \textbf{Configuración} & \textbf{AUROC} & \textbf{Nota} \\
\midrule
VICReg nativo & ConvNeXt-tiny, sin preentrenar, sin \emph{dropout} & 0.756 & control aleatorio real: 0.748 \\
VICReg + ImageNet & sin \emph{dropout} & sin evaluar & sobreajuste detectado antes del \emph{probe} \\
VICReg + ImageNet + \emph{dropout} 0.3 & \emph{dropout}=0.3 & abandonado & colapso progresivo, no regularización \\
\textbf{VICReg + ImageNet + \emph{dropout} 0.1} & \emph{dropout}=0.1 & \textbf{0.829} & mejor corrida de la sesión \\
BYOL (ResNet-18 preentrenado) & --- & \textbf{0.822} & mejora consistente (bootstrap $P{=}0.991$) \\
DINOv2 (ViT-S \emph{patch} 60) & --- & 0.711 & apenas sobre el azar ($P{=}0.659$), consistente con la campaña oficial (0.717) \\
\bottomrule
\end{tabular}
\caption{Resultados exploratorios de la sesión del 20-09-2026, mismo split
que la campaña oficial (\texttt{campaign\_full50}), 30 épocas. Cifras
tomadas de \texttt{SESSION\_HANDOFF.md}; no persistidas todavía como
\texttt{results.json} ni como figuras versionadas.}
\label{tab:extension}
\end{table}

\subsection{Una advertencia importante sobre el control ``aleatorio''}
\label{sec:caveat}

\begin{warnbox}
El script \texttt{evaluate\_frozen\_probe.py} construye el ``encoder
aleatorio'' de control reconstruyendo el \emph{backbone} a partir de la
configuración guardada del propio \emph{checkpoint}
(\texttt{build\_backbone(cfg)}). Si esa configuración tiene
\texttt{model.backbone.pretrained=true} --- como ocurre en las corridas de
VICReg con ImageNet de esta sesión ---, el control ``aleatorio'' \textbf{no
es en realidad un encoder sin entrenar}: también carga los pesos de
ImageNet, solo que sin el ajuste SSL posterior. Por eso el $0.811$ de AUROC
que obtuvo ese control \emph{no es comparable} con los controles
verdaderamente aleatorios de las otras filas ($0.748$, $0.744$, $0.698$).
Esta corrección --- forzar siempre inicialización aleatoria real en el
control, sin importar \texttt{pretrained} --- quedó identificada pero
\textbf{sin implementar} al cierre de la sesión.
\end{warnbox}

Comparando en cambio contra el control verdaderamente aleatorio de la
campaña oficial (0.748), VICReg con ImageNet + \emph{dropout} 0.1 mejora
$+0.081$, en línea con la mejora de BYOL ($+0.078$ frente a su propio
control real, 0.744). Pero el propio VICReg con ImageNet, comparado contra
\emph{su análogo} con ImageNet sin entrenamiento SSL, solo mejora $+0.018$
con evidencia estadística débil (\emph{bootstrap} $P=0.71$): es decir, gran
parte de la ganancia observada podría deberse al preentrenamiento en
ImageNet, no al entrenamiento SSL sobre OASIS3 en sí. Lo mismo aplica, sin
cuantificar todavía, a la mejora de BYOL con ResNet-18.

\begin{tldrbox}
Con \emph{backbones} preentrenados en ImageNet, tanto VICReg como BYOL
llegan a AUROC $\approx0.82$--$0.83$ --- notablemente mejor que la campaña
oficial sin preentrenar ($0.78$ y $0.74$). Pero todavía no se sabe cuánto de
esa mejora viene del entrenamiento SSL sobre OASIS3 y cuánto ya venía
``de fábrica'' con ImageNet, porque el control usado para aislar esa
pregunta tiene el error descrito arriba. DINOv2 se mantiene marginal en
ambas sesiones.
\end{tldrbox}

% ============================================================
\section{Limitaciones, lecciones y próximos pasos}
\label{sec:limits}
% ============================================================

\begin{itemize}
\item \textbf{Un solo split, una sola semilla.} Todos los intervalos de
  confianza de este informe son \emph{bootstrap} condicionado a un único
  split de sujetos; no reflejan variabilidad entre splits ni entre semillas
  de inicialización.
\item \textbf{Test pequeño para el efecto que se busca.} 134 sujetos (31
  AD) es una muestra modesta para diferencias de AUROC de una o dos
  centésimas; los intervalos de confianza de la campaña oficial son, en
  consecuencia, anchos.
\item \textbf{Colapso parcial en BYOL y DINOv2.} El rango efectivo bajo
  (Tabla~\ref{tab:rank}) sugiere que, antes de correr más épocas, vale la
  pena revisar hiperparámetros específicos anti-colapso de cada método
  (por ejemplo, la temperatura del profesor en DINOv2, o el \emph{momentum}
  EMA de BYOL) más que simplemente extender el entrenamiento.
\item \textbf{Confusión preentrenamiento-ImageNet vs.\ SSL-OASIS3.} Como se
  detalla en la Sección~\ref{sec:caveat}, el control aleatorio de
  \texttt{evaluate\_frozen\_probe.py} debe forzarse a pesos aleatorios
  reales independientemente de \texttt{model.backbone.pretrained}, para
  poder atribuir correctamente la mejora observada.
\item \textbf{Nada de la sesión exploratoria está en control de versiones
  todavía.} Los \texttt{.json} del \emph{probe}, las figuras nuevas y los
  \emph{checkpoints} de esa sesión no están comiteados; persistirlos en
  \texttt{reports/frozen\_probe/} es el paso pendiente más simple y de
  mayor valor inmediato.
\item \textbf{t-SNE de embeddings}: existe para una corrida antigua de 270
  sujetos, pero no se ha regenerado para la cohorte completa de 981
  sujetos ni para las corridas de esta campaña.
\end{itemize}

% ============================================================
\section{Cómo reproducir estos resultados}
\label{sec:repro}
% ============================================================

Entorno: \texttt{oasis3-ssl-workbench/.venv/Scripts/python.exe} (Windows).
Comando base de la campaña oficial de 50 épocas (varía \texttt{model=} y
\texttt{loss=} para cada método, ver \texttt{scripts/run\_full50.sh}):

\begin{quote}
\begin{verbatim}
.venv/Scripts/python.exe -m t5_ssl.train \
  data=oasis3_full50 trainer=full50 trainer.accelerator=cpu \
  trainer.enable_progress_bar=false
\end{verbatim}
\end{quote}

Evaluación del \emph{probe} congelado sobre un experimento ya entrenado:

\begin{quote}
\begin{verbatim}
.venv/Scripts/python.exe scripts/evaluate_frozen_probe.py \
  logs/<experimento> --checkpoint best
\end{verbatim}
\end{quote}

\begin{ideabox}
En una máquina con poca memoria (esta se corrió con 7.7\,GB de RAM y hasta
tres entrenamientos concurrentes) es obligatorio reducir el caché de LMDB:
\texttt{data.train\_cache\_capacity=80 +data.eval\_cache\_capacity=16}. Con
el valor por defecto (1\,300) el proceso termina en \emph{out-of-memory}.
\end{ideabox}

% ============================================================
\appendix
\section{Glosario exprés}
% ============================================================

\begin{description}[leftmargin=3.2cm,style=nextline]
\item[AUROC] Área bajo la curva ROC; $0.5$ equivale a una clasificación al
  azar, $1.0$ a una separación perfecta entre AD y HC.
\item[Backbone / encoder] Red convolucional o \emph{transformer} que
  convierte una imagen en un vector de características; es la única parte
  que se conserva después del preentrenamiento SSL.
\item[EMA] \emph{Exponential Moving Average}: una copia de los pesos que se
  actualiza como promedio móvil de los pesos ``reales'', nunca por
  gradiente directo.
\item[\emph{Frozen probe}] Protocolo de evaluación que congela el encoder y
  entrena solo un clasificador lineal encima, usando las etiquetas
  clínicas únicamente en este paso.
\item[\emph{Rango efectivo} (\emph{effective rank})] Medida continua de
  cuántas direcciones del espacio de embeddings concentran la varianza;
  bajo $\Rightarrow$ posible colapso parcial.
\item[SSL] Aprendizaje autosupervisado (\emph{self-supervised learning}):
  entrenar sin etiquetas humanas, usando la propia estructura de los datos
  (aquí, ``dos vistas del mismo cerebro'') como señal de supervisión.
\item[\emph{Stop-gradient}] Operación que impide que el gradiente fluya
  hacia una rama de la red (p.\,ej.\ la rama \emph{target} de BYOL o la
  profesora de DINOv2), aunque sus salidas se usen para calcular la
  pérdida.
\end{description}

\end{document}
